
\begin{center}
\icmltitle{Supplement to ``Improving Generalization with Flat Hilbert Bayesian Inference''}
\end{center}

The Appendix will be structured as follows: Section \ref{appendix: proofs} presents the proofs to our theoretical results. Next, Section \ref{appendix: additional experiments} is dedicated for the additional experiments to assess the effectiveness of our method, including the effect of the number of particles, and the effect of kernel choices. Then, Section \ref{appendix: experimental details} discusses the related experimental details, including the chosen hyperparameters or data augmentations used in our experiments.

\section{All Proofs}
\label{appendix: proofs}

We introduce a few additional notations for the sake of the missing proofs of the main theoretical results. Given a RKHS $\gH$ equipped with the inner product $\left<\cdot, \cdot\right>_\gH$ and the norm operator $\|\cdot \|_\gH$. We define the single-sample loss function on the functional space $\gH$ to be a map:

\begin{align*}
    \Tilde{\ell}: \gH^d \times \gX \times \gY &\to \mathbb{R}\\
    (\vf, x, y) &\mapsto \Tilde{\ell}(\vf, (x, y)).
\end{align*}
 
Define the \textit{general functional loss} $\Tilde{L}_\gD(\vf) = \mathbb{E}_{(x, y) \sim \gD}[\Tilde{\ell}(\vf, (x, y))]$ and the \textit{empirical functional loss} $\Tilde{L}_\gS(\vf) = \sum_{i=1}^n\Tilde{\ell}(\vf, (x_i, y_i))$. Throughout the proof, we assume that the parameter space is bounded by $\|\vtheta\| \leq H$, the data is al bounded that $\|x\| \leq R, y \leq R$ for some $R, H \in \mathbb{R}$, and the loss function $\Tilde{l}$ is $s-$Hölder continuous. Indeed, this $s-$Hölder continuity property holds for a wide range of loss functions, including the MSE loss, Huber loss, Hinge loss, Smooth L1 loss, and Logistic loss.

We introduce the following lemmas that will be used throughout the proof of our main theorems.

\begin{lemma}[Approximation of RKHS functionals]
\label{lemma: RKHS approximation}

Let $d \in \mathbb{N}, \gX = [-R, R]^K$ for some $K \in \mathbb{R}$. Consider $\gK = \{f \in \gH: \|f\|_{\gH} \leq 1\}$ with $\gH$ induced by some Mercer kernel which is $\alpha$-Holder continuous for $\alpha \in (0, 1)$ with constant $C_K \geq 0$. Suppose $F$ is $s-$Holder continuous for $s \in (0, 1]$ with constant $C_f \geq 0$. There exists some $M_0 \in \mathbb{N}$ such that for every $M \in \mathbb{N}$ with $M > M_0$, by taking some fixed $\overline{t}= \{t_i\}$ with $N  \in \mathbb{N}$, we have a tanh neural network $\hat{G}$ with two hidden layers of widths at most $N(M-1)$ and $3\frac{N+1}{2}(5M)^N$ parameters satisfying 
    \begin{equation}
    \label{inequality: rkhs approximation}
        \sup_{f \in \gK}|F(f) - \hat{G}(f(\overline{t}))| \leq RC_F(\epsilon_K(\overline{t}))^s + \frac{7N^2R C_G}{M},
    \end{equation}
    with 
    \begin{equation*}
        C_G = C_F(1+\|K[\overline{t}]^{-1}\|_{op}\sqrt{N}C_K(h_{\overline{t}})^\alpha)^s,
    \end{equation*}
    where $K[\overline{t}]$ is the Gram matrix of $\overline{t}$.
\end{lemma}
\begin{proof}
    The proof can be found in \cite{rkhsapprox}
\end{proof}

\begin{lemma}[Product of RKHSs]
\label{lemma: product of rkhs}
    Given $n$ RKHSs $\gH_1, \gH_2, \cdots, \gH_n$, each defined on corresponding sets $X_1, X_2, \cdots, X_n$ with kernels $k_1(x_1, y_1), \cdots k_n(x_n, y_n)$ respectively. Then, $\gH = \bigotimes_{i = 1}^n = \gH_1 \times \gH_2 \times \cdots \times \gH_n$ is also an RKHS, with kernel $K$ that is the product of the individual kernels. 
\end{lemma}
\begin{proof}
    The product space $\gH = \bigotimes_{i=1}^n \gH_i$ consists of tuples of functions $(f_1, f_2, \cdots, f_n)$. Firstly, we define the inner product in $\gH$ as:
    \begin{equation*}
        \left< (f_1, f_2, \cdots, f_n), (g_1, g_2, \cdots, g_n)\right>_\gH = \sum^n_{i=1}\left<f_i, g_i\right>_{\gH_i}.
    \end{equation*}
    This definition naturally defines a Hilbert space structure on $\gH$ since each $\gH_i$ is a Hilbert space, and the sum of inner products is linear and positive definite. Now we define the kernel for the product space:
    \begin{equation*}
        k((x_1, x_2, \cdots, x_n), (y_1, y_2, \cdots, y_n)) = \prod_{i=1}^n k(x_i, y_i).
    \end{equation*}
    Notice that the pointwise product of positive definite kernels is a positive definite kernel, hence this kernel is valid.

    We now verify the reproducing property of $\gH$. Consider a function $f = (f_1, f_2, \cdots, f_n) \in \gH$, and evaluate the function at a point $(x_1, x_2, \cdots, x_n) \in \bigotimes_{i=1}^n X_i$. 

    The reproducing property in each individual RKHS $\gH_i$ implies that:
    \begin{equation*}
        f_i(x_i) = \left< f_i, k_i(x_i, \cdot) \right>_{\gH_i}.
    \end{equation*}
    Hence, for the function $f = (f_1, \cdots, f_n)$, we get:
    \begin{align*}
        f((x_1, x_2, \cdots, x_n)) &= (f_1(x_1), f_2(x_2),\cdots,f_n(x_n))\\
        &=(\left<f_1, k_1(x_1, \cdot)\right>_{\gH_1}, \left<f_2, k_2(x_2, \cdot)\right>_{\gH_2},\cdots,\left<f_n, k_n(x_n, \cdot)\right>_{\gH_n})\\
        &=\left<(f_1,f_2,\cdots,f_n),(k_1(x_1,\cdot),k_2(x_2,\cdot),\cdots k_n(x_n,\cdot))\right>_\gH.
    \end{align*}
    Thus, the reproducing property holds for the product space $\gH$. Since $\gH$ is a Hilbert space and the kernel $k$ satisfies the reproducing property, we conclude that $\gH = \bigotimes^n_{i=1} \gH_i$ is another RKHS.  
\end{proof}
\subsection{Proof of Proposition \ref{proposition 1}}
\label{proof: proposition 1}
\begin{proposition}
    Consider the problem of finding the distribution $\mathbb{Q}$ that solves:
    \begin{equation}
        \mathbb{Q}^* = \min_{\mathbb{Q} \ll \mathbb{P}_{\mathbf{\theta}}} \Bigg\{\mathbb{E}_{\theta \sim \mathbb{Q}}[\mathcal{L}_\mathcal{D}(\vtheta)] + \KL(\mathbb{Q} \| \mathbb{P}_{\mathbf{\theta}}) \Bigg\}
    \end{equation}
    where we search over $\mathbb{Q}$ absolutely continuos w.r.t $\mathbb{P}_\vtheta$, and the second term is the regularization term. The closed-form solution to this problem is the \textbf{population posterior} whose density has the form:
    \begin{equation*}
        q^*(\vtheta) \propto \exp(-\mathcal{L}_\mathcal{D}(\vtheta))p(\vtheta).
    \end{equation*}
\end{proposition}
\begin{proof}
    This proposition is the general case of \textbf{Theorem 3.1} by \cite{va2023_flat_Bayes}. Denote $q(\cdot)$ as the density function of $\mathbb{Q}$. We have: \begin{align*}
        \mathbb{E}_{\theta \sim \mathbb{Q}}[\mathcal{L}_\mathcal{D}(\vtheta)] + \KL(\mathbb{Q} \| \mathbb{P}_{\mathbf{\theta}}) = \int_\Theta \mathcal{L}_\mathcal{D}(\vtheta)q(\vtheta) d\vtheta + \int_\Theta q(\vtheta) \log\frac{q(\vtheta)}{p(\vtheta)}d\vtheta.
    \end{align*}
    The Lagrangian is given by:
    \begin{equation*}
        L(q, \alpha) = \int_\Theta \mathcal{L}_\mathcal{D}(\vtheta)q(\vtheta) d\vtheta + \int_\Theta q(\vtheta) \log\frac{q(\vtheta)}{p(\vtheta)}d\vtheta + \alpha(\int q(\vtheta)d\vtheta - 1).
    \end{equation*}
    Taking derivative with respect to $q(\vtheta)$, it follows
    \begin{align*}
        &\mathcal{L}_\mathcal{D} + \log q(\vtheta) + 1 - \log p(\vtheta) + \alpha = 0,\\
        &q(\vtheta) = \exp(-\mathcal{L}_\mathcal{D}(\theta))p(\vtheta)\exp(-\alpha -1),\\
    \end{align*}
    which implies that
    \begin{equation*}
        q(\vtheta) \propto \exp(-\mathcal{L}_\mathcal{D}(\vtheta))p(\vtheta).
    \end{equation*}
    Then, the optimal solution is the population posterior $p(\vtheta|\mathcal{S})$, which concludes the proof.
\end{proof}
\subsection{Proof of Theorem \ref{theorem: flow sharpness}}
\label{proof: flow sharpness}

\begin{theorem} 
For any $\rho > 0$ and any distribution $\gD$, with probability $1 - \delta$ over the choice of the training set $\gS \sim \gD^n$, 
\begin{align*}
     \Tilde{L}_\gD (\vf) &\leq \max_{\|\vf' - \vf\|_{\gH^d} \leq \rho}\Tilde{L}_\gS(\vf') + \\
     % RC_F(\epsilon_K(\overline{t}))^s + \frac{7N^2R C_G}{M}\\
         &+ \sqrt{\frac{N' \log \Bigg(1 + \frac{C}{\rho^2 P^2}\Bigg(1+ \sqrt{\frac{\log(N)}{N'}}\Bigg)^2 \Bigg)+4\log\frac{n}{\delta}+8\log(6n+3k)}{n-1}}.\\
\end{align*}
\end{theorem}

\begin{proof}

$\Tilde{\ell}$ is a functional that maps from $\gH^d \times \gX \times \gY$ to $\mathbb{R}$. Notice that $\gH^d$ is a RKHS, $\gX = \mathbb{R}^a$ and $\gY = \mathbb{R}^b$ for some $a, b \in \mathbb{Z}$ are Euclidean spaces, which are also instances of RKHS. Moreover, the product of RKHS's is also a RKHS according to Lemma \ref{lemma: product of rkhs}. Hence, $\gH^d \times \gX \times \gY$ is also a RKHS. According to Lemma \ref{lemma: RKHS approximation}, there exists $N$ points $\overline{\vtheta} = [\vtheta_i]^N_{i = 1} \subset \Theta$, and a two-layer neural network $G_\mW$ parameterized by $\mW$ so that
\begin{equation*}
    |\Tilde{\ell}(\vf, x, y) - G_\mW(\vf(\overline{\vtheta}), x, y)| \leq  RC_F(\epsilon_K(\overline{t}))^s + \frac{7N^2R C_G}{M},
    \end{equation*}
    for every $(\vf, x, y) \in \gH^d \times \gX \times \gY$.  Consider $\vf' \in \gH^d$ so that $\|\vf' - \vf\| \leq \rho$. Recall that we have $\overline{\vtheta} = [\vtheta_i]_{i=1}^N$ where each $\vtheta_i \in \mathbb{R}^d$, and $f(\overline{\vtheta}) = [f(\vtheta_1), \cdots, f(\vtheta_N)] \in \mathbb{R}^{N'}$ where $N' =Nd$. Then, we have:
    \begin{equation*}
        \|f(\overline{\vtheta}) - f'(\overline{\vtheta})\|^2_2 = \Big(\sum \|f(\vtheta_i) - f'(\vtheta_i)\|\Big)^{1/2}.
    \end{equation*}
    Now let $f= [f_1, f_2, \cdots, f_d]$ and $f' = [f_1', f_2', \cdots, f_d']$ where $f_i, f'_j \in \gH$, we have:
    \begin{equation*}
        \|f(\vtheta_i) - f'(\vtheta_i)\| = \sum^d_{j=1} |f_j(\vtheta_i) - f'_j(\vtheta_i)| = \sum^d_{j=1}|\left<k(\vtheta_i, \cdot), f_j - f'_j \right>| \leq \sum^d_{j=1}\|k(\vtheta_i, \cdot)\| \|f_j - f'_j\| = k(\vtheta_i, \vtheta_i)\|f - f'\|.
    \end{equation*}
    Therefore, it follows:
    \begin{equation*}
        \|f(\overline{\vtheta}) - f'(\overline{\vtheta})\| \leq \Big( \sum k(\vtheta_i, \vtheta_i)\Big)^{1/2} \|f - f'\|.
    \end{equation*}
    Under the assumption that $k(\vtheta, \vtheta) \leq C$ (this is true for the widely used kernels, for example, $C = 1$ for the RBF kernel) it implies $   |\vf(\overline{\vtheta}) - \vf'(\overline{\vtheta})| \leq P \| \vf - \vf'\|_{\gH^d} \leq P\rho$ where $P = CN^{1/2}$. Denote $\Tilde{\vtheta} = \vf(\overline{\vtheta}) \in \mathbb{R}^{N'}$ for some $N' \in \mathbb{Z}$, by invoking the inequality from \cite{sam}, let $\rho' = \rho P$, it follows that: 
    
    % Let $L_\vtheta$ be the evaluation functional over $\gH^d$. According to the definition of RKHS, for all $\vtheta \in \Theta$, $L_\vtheta$ is continuous at every $\vf \in \gH^d$. Equivalently, $L_\vtheta$ is a bounded operator on $\gH^d$, which means there exists $M_\vtheta > 0$ such that $|L_\vtheta (\vf)| = |\vf(\vtheta)| \leq M \|\vf\|_{\gH^d}$. 
    

    %One have that:

    % \begin{align*}
    %     \Big|G_\mW(\vf'(\overline{\vtheta})) - G_\mW(\vf'(\overline{\vtheta}))\Big|& =  \Big| W_1 \tanh (W_0^\top [\vf'(\overline{\vtheta}), x, y] + b_0) - W_1 \tanh (W_0 [\vf(\overline{\vtheta}), x, y]^\top + b_0)\Big|\\
    %     &\leq \| W_1\|\Big|\tanh(W_0[\vf(\overline{\vtheta}), x, y]^\top) - \tanh(W_0[\vf'(\overline{\vtheta}), x, y]^\top)\Big|\\
    %     &\leq \|W_1\| \|W_0([\vf(\overline{\vtheta}), x, y] - [\vf'(\overline{\vtheta}), x, y])\|\\
    %     &\leq \|W_1\| \|W_0\| \|\vf(\overline{\vtheta}) - \vf'(\overline{\vtheta})\| \leq TP \rho
    % \end{align*}
%for some constant $T$ that bounds the parameters of $G$. 



\begin{align*}
    \Tilde{L}_\gD (\vf) &= \mathbb{E}_{(x, y) \sim \gD}[\Tilde{\ell}(\vf, x, y)] \leq \mathbb{E}_{(x, y)\sim \gD}[G_\mW(\vf(\overline{\vtheta}), x, y)] +  RC_F(\epsilon_K(\overline{t}))^s + \frac{7N^2R C_G}{M}\\
    & = \mathbb{E}_{(x, y)\sim \gD}[G_\mW(\Tilde{\vtheta}, x, y)] +  RC_F(\epsilon_K(\overline{t}))^s + \frac{7N^2R C_G}{M}\\
    &\leq \max_{\|\Tilde{\vtheta}' - \Tilde{\vtheta}\|^2_2 \leq \rho'} \frac{1}{n}\sum_{i = 1}^n G_\mW(\Tilde{\vtheta}', x, y)  + h(M, N)\\
    &+ \sqrt{\frac{N' \log \Bigg(1 + \frac{C}{\rho'^2}\Bigg(1+ \sqrt{\frac{\log(N)}{N'}}\Bigg)^2 \Bigg)+4\log\frac{n}{\delta}+8\log(6n+3k)}{n-1}},
\end{align*}

where we defined $h(M, N) =  RC_F(\epsilon_K(\overline{t}))^s + \frac{7N^2R C_G}{M}$. By definition, a RKHS is a closed Hilbert space. Then, there exists a sequence $\{\vf'_n\}$ so that $\vf'_n(\overline{\vtheta})$ that gets arbitrarily close to $\Tilde{\vtheta}'$. Then, for any $\epsilon > 0$, it follows:

\begin{align*}
     \Tilde{L}_\gD (\vf) &\leq \max_{\|\Tilde{\vtheta}' - \Tilde{\vtheta}\|^2_2 \leq \rho'} \frac{1}{n}\sum_{i = 1}^n G_\mW(\Tilde{\vtheta}', x, y)  + h(M, N)\\
    &+ \sqrt{\frac{N' \log \Bigg(1 + \frac{C}{\rho'^2}\Bigg(1+ \sqrt{\frac{\log(N)}{N'}}\Bigg)^2 \Bigg)+4\log\frac{n}{\delta}+8\log(6n+3k)}{n-1}}\\
    & \leq \max_{\|\vf'(\overline{\vtheta}) - \vf(\overline{\vtheta})\|^2_2 \leq \rho P} \frac{1}{n}\sum_{i = 1}^n G_\mW(\vf'(\overline{\vtheta}), x, y)  + h(M, N) + \epsilon \gO(1)\\
    &+ \sqrt{\frac{N' \log \Bigg(1 + \frac{C}{\rho^2P^2}\Bigg(1+ \sqrt{\frac{\log(N)}{N'}}\Bigg)^2 \Bigg)+4\log\frac{n}{\delta}+8\log(6n+3k)}{n-1}}\\
    &\leq \max_{\|\vf' - \vf\|^2_2 \leq \rho} \frac{1}{n}\sum_{i = 1}^n G_\mW(\vf'(\vtheta), x, y)  + h(M, N) + \epsilon \gO(1)\\
    &+ \sqrt{\frac{N' \log \Bigg(1 + \frac{C}{\rho^2 P^2}\Bigg(1+ \sqrt{\frac{\log(N)}{N'}}\Bigg)^2 \Bigg)+4\log\frac{n}{\delta}+8\log(6n+3k)}{n-1}}\\
    &\leq \max_{\|\vf' - \vf\|^2_2 \leq \rho}\Tilde{L}_\gS(\vf') + h(M,N) + \epsilon \gO(1)\\
    &+ \sqrt{\frac{N' \log \Bigg(1 + \frac{C}{\rho^2 P^2}\Bigg(1+ \sqrt{\frac{\log(N)}{N'}}\Bigg)^2 \Bigg)+4\log\frac{n}{\delta}+8\log(6n+3k)}{n-1}}.\\
\end{align*}
This is true for any $\epsilon > 0$.  Moreover, we can choose $\epsilon_K$ and $M$ to be arbitrarily small so that $h(M, N) \to 0$. Hence, it implies
\begin{align*}
     \Tilde{L}_\gD (\vf) &\leq \max_{\|\vf' - \vf\|^2_2 \leq \rho}\Tilde{L}_\gS(\vf') + \\
         &+ \sqrt{\frac{N' \log \Bigg(1 + \frac{C}{\rho^2 P^2}\Bigg(1+ \sqrt{\frac{\log(N)}{N'}}\Bigg)^2 \Bigg)+4\log\frac{n}{\delta}+8\log(6n+3k)}{n-1}},\\
\end{align*}
which concludes our proof.
\end{proof}
\subsection{Proof of Theorem \ref{theorem: pac-bayes}}
Now we can prove the Theorem \ref{theorem: pac-bayes}. We restate the theorem

\label{proof: pac-bayes}
\begin{theorem}
    For any target distribution $p$, reference distribution $q$, and any $\rho > 0$, we have the following bound between the general KL loss and the empirical KL loss
\begin{align*}
    \KL\Bigg(q_{[\vf]} || p\Big(\vtheta|\gD\Big)\Bigg) &\leq \max_{\vf' \in \gB_\rho(\vf)}\KL\Bigg(q_{[\vf']} || p\Big(\vtheta | \gS\Big)\Bigg)\\
         &+ \sqrt{\frac{N' \log \Bigg(1 + \frac{C}{\rho^2 P^2}\Bigg(1+ \sqrt{\frac{\log(N)}{N'}}\Bigg)^2 \Bigg)+4\log\frac{n}{\delta}+8\log(6n+3k)}{n-1}}.\\
    \end{align*}
\end{theorem}

\begin{proof}

Consider the left-hand side, we have:

\begin{align*}
    \KL(q_{[\vf]} \| p(\vtheta|\gD)) &= \int q_{[\vf]}(\vtheta)\Bigg(\gL_\gD(\vtheta)+\log \frac{q_{[\vf]}(\vtheta)}{p(\vtheta)} + \log Z_\gD \Bigg)d\vtheta\\
    &=\int q_{[\vf]}(\vtheta)\Bigg(\mathbb{E}_{(x, y) \sim \gD} \ell(\vtheta, x, y)+\log \frac{q_{[\vf]}(\vtheta)}{p(\vtheta)} + \log Z_\gD\Bigg)d\vtheta\\
    &=\mathbb{E}_{(x, y)\sim \gD}\Bigg[\int q_{[\vf]}(\vtheta)\Big(\ell(\theta; x, y)+\log\frac{q_{[\vf]}(\vtheta)}{p(\vtheta)}\Big) d\vtheta \Bigg] + \int q_{[\vf]}(\vtheta) \log Z_\gD d\vtheta,
\end{align*}

where $Z_\gD$ is the normalizing constant.  On the other hand, we also have:

\begin{align*}
    \KL(q_{[\vf]} \| p(\vtheta|\gS)) &= \int q_{[\vf]}\Bigg((\vtheta)\gL_\gS(\vtheta)+\log \frac{q_{[\vf]}(\vtheta)}{p(\vtheta)}+\log Z_\gS\Bigg) d\vtheta\\
    &=\int q_{[\vf]}(\vtheta)\Bigg(\frac{1}{n}\sum_{i = 1}^n \ell(\vtheta, x_i, y_i)+\log \frac{q_{[\vf]}(\vtheta)}{p(\vtheta)}+\log Z_\gS\Bigg)d\vtheta\\
    &=\frac{1}{n}\sum_{i=1}^n\Bigg[\int q_{[\vf]}(\vtheta)\Big(\ell(\theta; x_i, y_i)+\log\frac{q_{[\vf]}(\vtheta)}{p(\vtheta)} + \log Z_\gS\Big) d\vtheta \Bigg] + \int q_{[\vf]}(\vtheta) \log Z_\gS d\vtheta,
\end{align*}

where $Z_\gS$ is the normalizing constant. We define $\Tilde{L}$ to be the functional such that:
\begin{align*}
    \Tilde{L}: \gH^d \times \gX \times \gY &\to \mathbb{R}\\
    (\vf, x, y) &\mapsto \Tilde{L}(\vf, x, y) = \int q_{[\vf]}(\vtheta)\Big(\ell(\theta; x, y)+\log\frac{q_{[\vf]}(\vtheta)}{p(\vtheta)}\Big) d\vtheta.
\end{align*}

According to Theorem \ref{theorem: flow sharpness}, we have:
\begin{align}
        \Tilde{L}_\gD(\vf) &\leq \max_{\|\vf' - \vf\|_{\gH^d} \leq \rho} \Tilde{L}_\gS(\vf')\\
        \label{eq: ineq}
         &+ \sqrt{\frac{N' \log \Bigg(1 + \frac{C}{\rho^2 P^2}\Bigg(1+ \sqrt{\frac{\log(N)}{N'}}\Bigg)^2 \Bigg)+4\log\frac{n}{\delta}+8\log(6n+3k)}{n-1}}.
\end{align}

Moreover, notice that we have $Z_\gD = \mathbb{E}[\exp(\gL_\gD(\vtheta))p(\vtheta)]$, $Z_\gS = \mathbb{E}[\exp(\gL_\gS(\vtheta))p(\vtheta)]$. Let us denote $\mathrm{VCdim}(\mathcal{F})=d_{\mathrm{VC}}$ where $\mathcal{F} = \{f_\vtheta: \vtheta \in \Theta\}$. According to Theorem 6.11 in \cite{uml} (Page 51), with the probability at least $1 - \delta$, we have 
\[
\sup_{\vtheta\in\Theta}|\mathcal{L}_{D}\left(\vtheta\right)-\mathcal{L}_{S}\left(\vtheta\right)|\leq\frac{4+\sqrt{\log\tau_{\mathcal{H}}\left(2n\right)}}{\delta\sqrt{2n}}.
\]
Note that the growth function $\tau_{\mathcal{H}}\left(2n\right)\leq\left(\frac{2en}{d_{\mathrm{VC}}}\right)^{d_{\mathrm{VC}}}$, we obtain
\[
\sup_{\vtheta\in\Theta}|\mathcal{L}_{D}\left(\vtheta\right)-\mathcal{L}_{S}\left(\vtheta\right)|\leq\frac{4+\sqrt{d_{\mathrm{VC}}\log\frac{2en}{d_{\mathrm{VC}}}}}{\delta\sqrt{2n}}.
\]
Therefore, for all $\vtheta \in \Theta$, we gain
\[
\exp\left(\mathcal{L}_{D}\left(\vtheta\right)\right)\leq\exp\left(\mathcal{L}_{S}\left(\vtheta\right)\right)\exp\left(\frac{4+\sqrt{d_{\mathrm{VC}}\log\frac{2en}{d_{\mathrm{VC}}}}}{\delta\sqrt{2n}}\right).
\]
It follows that
\[
\int\exp\left(\mathcal{L}_{D}\left(\vtheta\right)\right)p\left(\vtheta\right)d\vtheta\leq\exp\left(\frac{4+\sqrt{d_{\mathrm{VC}}\log\frac{2en}{d_{\mathrm{VC}}}}}{\delta\sqrt{2n}}\right)\int\exp\left(\mathcal{L}_{S}\left(\vtheta\right)\right)p\left(\vtheta\right)d\vtheta.
\]
Which is equivalent to
\[
Z_{D}\leq\exp\left(\frac{4+\sqrt{d_{\mathrm{VC}}\log\frac{2en}{d_{\mathrm{VC}}}}}{\delta\sqrt{2n}}\right)Z_{S}.
\]
Therefore, we have
\begin{equation*}
    \label{ineq: zs_zd}
    \log Z_{D}\leq\log Z_{S}+\frac{4+\sqrt{d_{\mathrm{VC}}\log\frac{2en}{d_{\mathrm{VC}}}}}{\delta\sqrt{2n}}
\end{equation*}


% \begin{equation}
% \label{eq: ineq Z}
%     \log Z_\gD \leq \log Z_\gS + \frac{K}{\sqrt{n}}
% \end{equation}

% almost surely with some constant $K > 0$.
% \begin{align*}
%     \log Z_\gS &= \int \exp(\gL_\gS(\vtheta))p(\vtheta)d\vtheta \\
%     &\geq \int \log(\exp(\gL_\gS(\vtheta))p(\vtheta))d\vtheta \\
%     & = \int \gL_\gS(\vtheta)+\log(p(\vtheta))d\vtheta\\
%     & \geq \int (\gL_\gD(\vtheta) - C(\frac{1}{\sqrt{n}}))+\log(p(\vtheta))d\vtheta \\
%     & = \int \log(\exp(\gL_\gD(\vtheta))p(\vtheta))d\vtheta - \int C(\frac{1}{\sqrt{n}})d\vtheta\\
%     & = \int \log(\exp(\gL_\gD(\vtheta))p(\vtheta))d\vtheta - \int C(\frac{1}{\sqrt{n}})d\vtheta\\
%     & = \int \log(\exp(\gL_\gD(\vtheta))p(\vtheta))d\vtheta - M^d C(\frac{1}{\sqrt{n}})\\
% \end{align*}

% Moreover, the model and data spaces are bounded, so $\gL_\gD(\vtheta)$ and $\gL_\gS(\vtheta)$ are bounded. Then, there exists constants $d, D$ such that $d \leq \log Z_\gS, \log Z_\gD \leq D$, which also implies $d \leq \int q_{[\vf]}\log Z_\gS d\vtheta \leq D$ and $d \leq \int q_{[\vf]}\log Z_\gD d\vtheta \leq D$. It follows that for all $\vf, \vf' \in \gH^d$:

% \begin{equation}
% \label{eq: ineq Z}
%     \int q_{[\vf]}(\vtheta) \log Z_\gD d\vtheta \leq \int q_{[\vf']}(\vtheta) \log Z_\gS d\vtheta + D - d.
% \end{equation}

Combining the Inequalities \ref{eq: ineq} and \ref{ineq: zs_zd}, it follows that:
\begin{align*}
    \KL\Bigg(q_{[\vf]} || p\Big(\vtheta|\gD\Big)\Bigg) &\leq \max_{\|\vf' - \vf\|_{\gH^d} \leq \rho}\KL\Bigg(q_{[\vf']} || p\Big(\vtheta | \gS\Big)\Bigg)\\
         &+ \sqrt{\frac{N' \log \Bigg(1 + \frac{C}{\rho^2 P^2}\Bigg(1+ \sqrt{\frac{\log(N)}{N'}}\Bigg)^2 \Bigg)+4\log\frac{n}{\delta}+8\log(6n+3k)}{n-1}} + \frac{4+\sqrt{d_{\mathrm{VC}}\log\frac{2en}{d_{\mathrm{VC}}}}}{\delta\sqrt{2n}}.
\end{align*}
which concludes our proof.
\end{proof}
\section{Additional Experiments}
\label{appendix: additional experiments}
\subsection{Effect of \#particles}
\label{appendix: number of particles}
\begin{wrapfigure}{r}{0.4\textwidth}
   \vspace{-13mm}
    \includegraphics[width=0.9\linewidth]{figs/runtime.pdf}
    \caption{Runtime by $\textsc{\#particles}$.}
    \label{fig:wt_a}
\end{wrapfigure}

To understand the impact of varying the number of particles, we experiment on the seven Natural datasets, reporting both accuracy and per-epoch runtime. We compared FHBI with SVGD and SAM. Figure \ref{fig:wt_a} and Table \ref{table: accuracy by particles} indicate that multiple particles result in significant performance improvements compared to a single particle. However, while increasing the number of particles enhances performance, it introduces a tradeoff regarding runtime and memory required to store the models. Based on these observations, we found that using $\textsc{\#particles} = 4$ provides an optimal balance between performance gains and computational overhead.

\begin{table}[h]
    \centering
    \renewcommand\arraystretch{1.2}
    \begin{tabular}{
>{\columncolor[HTML]{FFFFFF}}l 
>{\columncolor[HTML]{FFFFFF}}c 
>{\columncolor[HTML]{FFFFFF}}c 
>{\columncolor[HTML]{FFFFFF}}c 
>{\columncolor[HTML]{FFFFFF}}c 
>{\columncolor[HTML]{FFFFFF}}c cc}
\hline
\textsc{\#Particles} &
  {CIFAR100} &
  {Caltech101} &
  {DTD} &
  {Flower102} &
  {Pets} &
  {SVHN} &
  {Sun397} \\ \hline
1p (SAM) &
  {\color[HTML]{1F2328} 72.7} &
  {\color[HTML]{1F2328} 90.3} &
  {\color[HTML]{1F2328} 71.4} &
  {\color[HTML]{1F2328} 99.0} &
  {\color[HTML]{1F2328} 90.2} &
  84.4 &
  52.4 \\ \hline
4p  & 74.8          & 93.0          & 74.3 & \textbf{99.4} & 92.4          & 87.5          & 56.5          \\ \hline
10p & \textbf{75.0} & \textbf{93.2} & \textbf{75.0} & 99.1          & \textbf{92.4} & \textbf{87.9} & \textbf{58.3} \\ \hline
\end{tabular}
\caption{Accuracy by $\textsc{\#particles}$.}
    \label{table: accuracy by particles}
\end{table}

\subsection{Effect of kernel choice}
\label{appendix: kernel choices}
The implementation of FHBI relies on the choice of the kernel $k$. In our experiments, we selected the RBF kernel due to its widespread use in the kernel methods literature, known for its strong representational capabilities and its ability to balance underfitting and overfitting through the kernel width parameter $\sigma$. To evaluate the impact of different kernel choices, we tested our method on the four Specialized datasets using the polynomial kernel of degree 10 as a comparison. The results, summarized in Table \ref{table: kernel choice}, indicate that while the polynomial kernel slightly underperforms relative to the RBF kernel, the difference is minimal, with a performance gap of less than 0.3\%.

% Please add the following required packages to your document preamble:
% \usepackage[table,xcdraw]{xcolor}
% Beamer presentation requires \usepackage{colortbl} instead of \usepackage[table,xcdraw]{xcolor}
\begin{table}[h]
\centering
\caption{Classification accuracy on the Specialized datasets with different kernel  choices}
\label{table: kernel choice}
\begin{tabular}{c|cccc|c}
{\color[HTML]{000000} Kernel}            & {\color[HTML]{000000} Camelyon}      & {\color[HTML]{000000} EuroSAT}       & {\color[HTML]{000000} Resisc45}      & {\color[HTML]{000000} Retinopathy}   & {\color[HTML]{000000} AVG}           \\ \hline
{\color[HTML]{000000} RBF}               & {\color[HTML]{000000} \textbf{85.3}} & {\color[HTML]{000000} \textbf{95.0}} & {\color[HTML]{000000} \textbf{87.2}} & {\color[HTML]{000000} \textbf{79.6}} & {\color[HTML]{000000} \textbf{86.8}} \\ \hline
{\color[HTML]{000000} Polynomial (d=10)} & {\color[HTML]{000000} 85.0}          & {\color[HTML]{000000} 94.9}          & {\color[HTML]{000000} 86.8}          & {\color[HTML]{000000} 79.2}          & {\color[HTML]{000000} 86.5}         
\end{tabular}
\end{table}

\subsection{Memory consumption and scalability to larger models}

To assess the scalability to larger models, we conduct an additional experiment on the larger ViT-B/16. In Tables \ref{tab:vram} and \ref{tab:trainingtime}, we report the memory usage and training time for the \texttt{CIFAR-100}, \texttt{Caltech-101}, and \texttt{Patch-Camelyon} datasets.


\begin{table}[h]
\centering
\caption{VRAM consumption (MB) on ViT-B/16 and ViT-L/16}
\label{tab:vram}
\begin{tabular}{c|ccc}
\hline
Architechture & CIFAR-100 & Caltech-101 & Patch-Camelyon \\ \hline
ViT-B/16               & 12541              & 12539                & 12535                   \\
ViT-L/16               & 33426              & 33684                & 32126                   \\ \hline
\end{tabular}
\end{table}

\begin{table}[h]
\centering
\caption{VRAM consumption (MB) on ViT-B/16 and ViT-L/16}
\label{tab:trainingtime}
\begin{tabular}{c|ccc}
\hline
Architechture & CIFAR-100 & Caltech-101 & Patch-Camelyon \\ \hline
ViT-B/16               & 5.55              & 5.66                & 5.46                   \\
ViT-L/16               & 17.45              & 17.35                & 17.02                   \\ \hline
\end{tabular}
\end{table}

\subsection{Sensitivity to the RBF kernel length scale}
As noted before, we initially tune the length scale parameter $\sigma$ from the candidates set $\{0.7, 1, 1.2\}$. To further investigate sensitivity to this hyperparameter, we expand the range to $\{0.1, 0.7, 1, 1.2, 2.5\}$, where we additionally include a small value of $\sigma = 0.1$ and a large value of $\sigma = 2.5$. The results on the \texttt{Natural} datasets, reported in Table \ref{tab:rbflengthsensitivity}, indicate that the performance remains robust within a reasonable range. We also observe a slight degradation when using extremely small values (e.g., $\sigma = 0.1$), where the model tends to overfit, or large values (e.g., $\sigma = 2.5$), where the model tends to underfit.

\begin{table}[h]
\centering
\caption{Sensitivity to the RBF length scale}
\label{tab:rbflengthsensitivity}
\begin{tabular}{c|ccccccc}
\hline
$\sigma$ & CIFAR-100     & Caltech-101   & DTD           & Flowers102    & Pets          & SVHN          & Sun397        \\ \hline
0.1                   & 72.1          & 91.6          & 74.0          & 97.9          & 90.2          & 85.3          & 52.2          \\
0.7                   & 73.8          & 91.8          & 73.3          & 98.7          & 92.4          & 86.7          & 56.1          \\
1                     & 73.6          & 92.7          & 72.7          & \textbf{99.1} & 91.9          & \textbf{87.3} & 54.3          \\
1.2                   & \textbf{74.1} & \textbf{93.0} & \textbf{74.3} & 98.3          & \textbf{92.4} & 86.4          & \textbf{56.5} \\
2.5                   & 69.2          & 90.9          & 69.4          & 97.5          & 90.9          & 84.6          & 52.6          \\ \hline
\end{tabular}
\end{table}

\section{Experimental Details}
\label{appendix: experimental details}
\subsection{Chosen Hyperparameters}
For each experiment, we conducted five runs of FHBI and reported the mean and standard deviation. All Bayesian methods were trained with four particles on the same set of LoRA parameters. We used ten warm-up epochs, batch size 64, the Gaussian kernel, and the cosine annealing learning rate scheduler for all settings. The experiments were run with PyTorch on a Tesla V100 GPU with 40GB of RAM. FHBI involves three hyperparameters: the learning rate $\epsilon$, ascent step size $\rho$, and kernel width $\sigma$. To tune these hyperparameters, we grid-search hyperparameters on the validation set, where the key hyperparameters are: the kernel width $\sigma$, the initial learning rate $\epsilon$, and the ascent step size $\rho$. The candidate sets are formed as $\epsilon \in \{0.15, 1, 1.5, 2.1, 2.5\}, \rho \in \{0.01, 0.03, 0.05\}, \sigma \in \{0.7, 1, 1.2\}$. The chosen hyperparameters are as follows $(\epsilon, \rho, \sigma)$: \texttt{CIFAR100} = $(0.15, 0.03, 1.2)$, \texttt{Caltech101} = $(2.1, 0.05, 1.2)$, \texttt{DTD} = $(0.15, 0.03, 1.2)$, \texttt{Flowers102} = $(0.15, 0.03, 1)$, \texttt{Pets} = $(0.15, 0.03, 1.2)$, \texttt{SVHN} = $(2.5, 0.01, 1)$, \texttt{Sun397} = $(0.15, 0.03, 1.2)$, \texttt{Patch-Camelyon} = $(2.1, 0.05, 1)$, \texttt{DMLab} = $(2.1, 0.03, 1)$, \texttt{EuroSAT} = $(2.5, 0.01, 1.2)$, \texttt{Resisc45} = $(1.5, 0.03, 1.2)$, \texttt{Diabetic- Retinopathy} = $(2.1, 0.03, 1)$, \texttt{Clevr-Count} = $(2.5, 0.01, 1)$, \texttt{Clevr-Dist} = $(1, 0.01, 1.2)$, \texttt{KITTI} = $(2.1, 0.05, 1)$, \texttt{dSprites-loc} = $(2.1, 0.05, 1)$, \texttt{dSprites-ori} = $(2.1, 0.03, 1.2)$, \texttt{smallNorb-azi} = $(1, 0.05, 1)$, \texttt{smallNorb-ele} = $(1, 0.03, 0.7)$.

\subsection{Data Augmentations}
Our implementation is based on the repository \href{https://github.com/synbol/Parameter-Efficient-Transfer-Learning-Benchmark/tree/main}{V-PETL}. Similar to this repository, we use a different data augmentation among the following three augmentations for each dataset. The data augmentations that we used for each setting are: 

\begin{itemize}
    \item For \texttt{CIFAR100, DTD, Flower102, Pets, Sun397}

\begin{lstlisting}[language=Python]
     self.transforms_train = transforms.Compose(
        [
            transforms.RandomResizedCrop(
                (self.size, self.size),
                scale=(self.min_scale, self.max_scale),
            ),
            transforms.RandomHorizontalFlip(self.flip_prob),
            transforms.TrivialAugmentWide()
            if self.use_trivial_aug
            else transforms.RandAugment(self.rand_aug_n, 
                                        self.rand_aug_m),
            transforms.ToTensor(),
            transforms.Normalize(mean=[0.485, 0.456, 0.406], 
                                std=[0.229, 0.224, 0.225])]),
            transforms.RandomErasing(p=self.erase_prob),
        ]
    )
    self.transforms_test = transforms.Compose(
        [
            transforms.Resize(
                (self.size, self.size),
            ),
            transforms.ToTensor(),
            transforms.Normalize(mean=[0.485, 0.456, 0.406], 
                                std=[0.229, 0.224, 0.225])]),
        ]
    )
\end{lstlisting}
\item For \texttt{Caltech101, Clevr-Dist, Dsprites-Loc, Dsprites-Ori, SmallNorb-Azi, SmallNorb-Ele}:

\begin{lstlisting}[language=Python]
    self.transform_train = transforms.Compose([
        transforms.Resize((224, 224)),
        transforms.ToTensor(),
        transforms.Normalize(mean=[0.485, 0.456, 0.406], 
                                std=[0.229, 0.224, 0.225])])
    self.transform_test = transforms.Compose([
        transforms.Resize((224, 224)),
        transforms.ToTensor(),
        transforms.Normalize(mean=[0.485, 0.456, 0.406], 
                                std=[0.229, 0.224, 0.225])])
\end{lstlisting}
\item For \texttt{Clevr-Count, DMLab, EuroSAT, KITTI, Patch Camelyon, Resisc45, SVHN, Diabetic Retinopathy}: 
\begin{lstlisting}
    from timm.data import create_transform
    self.transform_train = create_transform(
                input_size=(224, 224),
                is_training=True,
                color_jitter=0.4,
                auto_augment='rand-m9-mstd0.5-inc1',
                re_prob=0.0,
                re_mode='pixel',
                re_count=1,
                interpolation='bicubic',
            )
    aug_transform.transforms[0] = transforms.Resize((224, 224), 
                                                    interpolation=3)
    self.transform_test = transforms.Compose([
            transforms.Resize((224, 224)),
            transforms.ToTensor(),
            transforms.Normalize(mean=[0.485, 0.456, 0.406], 
                                std=[0.229, 0.224, 0.225])])
\end{lstlisting}
\end{itemize}

%\section{Additional Experiments}
%\input{sections/tables/runtime}